\section{Strategische Ma{\ss}nahmen}
\label{sec:strategie}

Dieser Teil berichtet und beurteilt nicht. Wertende Adjektive und
zukunftsgerichtete Konditionale bleiben ihm fern; die Beurteilung steht in
den Teilen~\ref{sec:kurs} bis \ref{sec:verdikt}. Wo das Unternehmen selbst
eine Erwartung ausspricht, wird sie als dessen Aussage gekennzeichnet und
nicht \"ubernommen.

\subsection{Island-Gold-Distrikt: Schachterweiterung und Ausbaustudie}

Die Schachterweiterung Phase~3+ am Island-Gold-Distrikt erreichte im Jahr
2025 eine Teufe von \num{\SchachtTiefe} Metern, nach Angabe des Unternehmens
\Pct{\SchachtAnteil} der Endteufe; die Fertigstellung ist f\"ur das Ende des Jahres 2026
vorgesehen.\quelleMDA{2025}{7}\quelleMerken{s9mdaxcacfxh}

Am \Datum{3.\,Februar 2026} wurde die Ausbaustudie f\"ur den Distrikt
ver\"offentlicht. Sie sieht gegen\"uber dem Basisplan vom Juni 2025 eine um
\Pct{\IGDReservenPlus} h\"ohere Reserve und eine Erweiterung der Magino-M\"uhle vor. Das
Unternehmen beziffert die Jahresproduktion der ersten zehn Jahre ab 2028 auf
\Unzen{\IGDProduktion} bei durchschnittlichen minenseitigen AISC von
\JeUnze{\IGDAISC} und den Kapitalwert nach Steuern bei einem Abzinsungssatz
von \Pct{\IGDZins} auf \MrdUSD{\IGDKapitalwert}; diese Angabe beruht auf einer
Preisannahme von \JeUnze{\IGDGoldpreis}.\quelleMDA{2025}{7}\quelleMerken{s10mdaxcacfxh}

\subsection{Lynn Lake und PDA}

F\"ur das Projekt Lynn Lake in Manitoba nennt das Unternehmen eine
Durchschnittsproduktion der ersten zehn Jahre von \Unzen{\LynnLakeProduktion}
bei minenseitigen AISC von \JeUnze{\LynnLakeAISC}; der Baubeginn ist f\"ur das
Fr\"uhjahr 2026 und die erste Produktion f\"ur das erste Halbjahr 2029
vorgesehen.\quelleMDA{2025}{13}\quelleMerken{s10mdaxcacfxbd} F\"ur das Projekt PDA
ist die erste Produktion zur Jahresmitte 2027
vorgesehen.\quelleErneut{s10mdaxcacfxbd}

\subsection{Reserven}

Die Reserven zum Jahresende 2025 betragen \MozReserven{\Reserven}, ein
Zuwachs von \Pct{\ReservenDelta} gegen\"uber dem Vorjahr, bei einem um
\Pct{\ReservenGehaltPlus} h\"oheren Gehalt von \num{\ReservenGehalt} Gramm je Tonne. Die
gemessenen und angezeigten Ressourcen betragen
\MozReserven{\ResourcenMI}.\quelleErneut{s10mdaxcacfxh} Es ist das
\num{\ReservenJahreFolge}.\ Jahr in Folge mit steigenden
Reserven.\quelleErneut{s10mdaxcacfxbd}

\subsection{Prognose gegen Ist}

Das Unternehmen berichtet f\"ur 2025 eine Produktion von
\Unzen{\ProduktionXXV} und stellt selbst fest, sie liege \emph{unter} der im
Jahresverlauf revidierten Prognose und um \Pct{\ProduktionRueckgangRund} unter dem Vorjahr; es
bezeichnet das Jahr aus betrieblicher Sicht als nicht kennzeichnend f\"ur die
eigene Ausf\"uhrungsbilanz.\quelleMDA{2025}{12}\quelleMerken{s11mdaxcacfxbc}

\begin{table}[htbp]
\centering
\caption{Prognose 2026 und Dreijahresausblick, beide vom Unternehmen
ver\"offentlicht. Die Kostenprognose f\"ur 2026 wurde im Zwischenbericht zum
30.\,Juni 2026 angehoben; beide St\"ande sind angegeben.}
\label{tab:prognose}
\begin{tabular}{lrr}
\toprule
Gr\"o{\ss}e & Prognose & Stand\\
\midrule
Produktion 2026 (Tsd.\ Unzen) & \TsdUnzen{\PrognoseProdUnten}--\TsdUnzen{\PrognoseProdOben} & Februar 2026\\
Produktion 2028 (Tsd.\ Unzen) & \TsdUnzen{\ProgAchtProdUnten}--\TsdUnzen{\ProgAchtProdOben} & Februar 2026\\
AISC 2026 & \JeUnze{\PrognoseAiscUnten}--\JeUnze{\PrognoseAiscOben} & Februar 2026\\
AISC 2026, angehoben & \JeUnze{\PrognoseAiscUntenNeu}--\JeUnze{\PrognoseAiscObenNeu} & Juli 2026\\
Gesamtinvestition 2026 (Mio.\,USD) & \MioUSDexakt{\PrognoseCapexUnten}--\MioUSDexakt{\PrognoseCapexOben} & Februar 2026\\
davon Wachstumskapital & \MioUSDexakt{\PrognoseWachstumsCapexUnten}--\MioUSDexakt{\PrognoseWachstumsCapexOben} & Februar 2026\\
davon Erhaltungskapital & \MioUSDexakt{\PrognoseErhaltCapexUnten}--\MioUSDexakt{\PrognoseErhaltCapexOben} & Februar 2026\\
\bottomrule
\end{tabular}
\end{table}

\FloatBarrier

Der Dreijahresausblick vom \Datum{4.\,Februar 2026} sieht eine Zunahme der
Produktion um \Pct{\ProgSechsProdDelta} im Jahr 2026 und um
\Pct{\ProgAchtProdDelta} bis 2028 sowie einen R\"uckgang der AISC um
\Pct{\ProgAchtAISCRueckgang} bis 2028 gegen\"uber 2025
vor.\quelleMDA{2025}{7}\quelleMerken{s11mdaxcacfxh}

Im Zwischenbericht zum \Datum{30.\,Juni 2026} wurde die Kostenprognose f\"ur
2026 angehoben: die Barkosten auf
\JeUnze{\PrognoseCashUnten}--\JeUnze{\PrognoseCashOben} und die AISC auf
\JeUnze{\PrognoseAiscUntenNeu}--\JeUnze{\PrognoseAiscObenNeu}.\quelleZB{6}\quelleMerken{s11zbxg}

\subsection{Wirkung im ersten Halbjahr 2026}

Der Umsatz des ersten Halbjahres 2026 betr\"agt \MioUSDexakt{\HJUmsatz} nach
\MioUSDexakt{\HJUmsatzVJ}, ein Zuwachs von \Pct{\HJUmsatzDelta}. Die
Produktion ging dabei um \PctBetrag{\HJProduktionDelta} auf
\Unzen{\HJProduktion} zur\"uck, der erzielte Goldpreis stieg um
\Pct{\HJGoldpreisDelta} auf \JeUnze{\HJGoldpreis}, die AISC stiegen um
\Pct{\HJAISCDelta} auf \JeUnze{\HJAISC}.\quelleZB{4}\quelleMerken{s11zbxe} Die Marge je Unze nach
Gleichung~\eqref{eq:marge-je-unze} stieg von \JeUnze{\HJMargeJeUnzeVJ} auf
\JeUnze{\HJMargeJeUnze}, also um \Pct{\HJMargeJeUnzeDelta}.

Der freie Cash-Flow des Halbjahres betr\"agt \MioUSDexakt{\HJFreierCashflow}
nach \MioUSDexakt{\HJFreierCashflowVJ}.\quelleErneut{s11zbxe} Die
Nettoliquidit\"at zum \Datum{30.\,Juni 2026} betr\"agt
\MioUSDexakt{\HJNettoliquiditaet}.\quelleZBA{2}\quelleMerken{s11zbaxc}
